\chapter{完整性约束、触发器与检查机制}

\section{完整性约束设计目标}

本章依据 \code{opengauss_setup/sql/init.sql}、\code{migrate_triggers_constraints_20260608.sql} 和 \code{test_triggers_constraints.sql}，说明当前系统已经落地的数据库完整性机制。完整性设计目标包括：实体可唯一标识、外键引用合法、业务取值范围合法、角色档案与账户角色一致、先修课程关系无自指和环路、选课过程不超容量且不发生时间冲突、成绩修改可以审计追踪。

与旧版本通过触发器维护 \code{selected_count} 的思路不同，当前系统在 3NF 迁移后删除了 \code{selected_count} 派生字段。本系统不再通过触发器同步该字段，而是通过 \code{enrollment} 表和视图统计当前已选人数。触发器的职责转向跨表业务完整性检查和审计记录维护，例如选课容量检查、时间冲突检查、先修课程合法性检查以及成绩变更日志写入。

\section{实体完整性：主键约束}

实体完整性主要通过主键实现。当前 \code{init.sql} 为所有核心表设置主键，例如 \code{department.dept_id}、\code{student.student_id}、\code{teacher.teacher_id}、\code{course.course_id}、\code{course_offering.offering_id}、\code{course_schedule.schedule_id}、\code{enrollment.enrollment_id}。这些主键保证每个实体或业务事件可以被唯一标识。

\section{参照完整性：外键约束}

参照完整性通过外键约束实现。典型外键包括：\code{major.dept_id} 引用 \code{department.dept_id}，\code{student.user_id} 引用 \code{user_account.user_id}，\code{course_offering.course_id} 引用 \code{course.course_id}，\code{course_schedule.offering_id} 引用 \code{course_offering.offering_id}，\code{enrollment.student_id} 引用 \code{student.student_id}，\code{score_change_log.enrollment_id} 引用 \code{enrollment.enrollment_id}。其中 \code{user_session.user_id} 和 \code{course_schedule.offering_id} 使用 \code{ON DELETE CASCADE}，用于自动清理会话和教学班时间片。

\section{用户自定义完整性：UNIQUE、CHECK、NOT NULL、DEFAULT}

本轮新增的 \code{migrate_triggers_constraints_20260608.sql} 在原有主键、外键、NOT NULL、DEFAULT 和 CHECK 约束基础上，进一步补充了业务唯一约束和合法性检查。表 \ref{tab:constraints-new} 总结了主要约束。

\begin{longtable}{p{0.23\textwidth}p{0.35\textwidth}p{0.33\textwidth}}
\caption{主要完整性约束}\label{tab:constraints-new}\\
\toprule
表 & 约束 & 业务含义 \\
\midrule
\endfirsthead
\toprule
表 & 约束 & 业务含义 \\
\midrule
\endhead
\bottomrule
\endfoot
\code{department} & \code{uq_department_name} & 学院名称唯一，避免同一学院重复建档。 \\
\code{major} & \code{uq_major_dept_name} & 同一学院内专业名称唯一。 \\
\code{classroom} & \code{uq_classroom_location} & 同一楼栋同一房间号只能对应一个教室实体。 \\
\code{course_prerequisite} & \code{chk_prereq_not_self} & 课程不能把自己作为自己的先修课程。 \\
\code{semester} & \code{chk_semester_date_range} & 学期起止日期、选课窗口起止时间必须合法。 \\
\code{user_session} & \code{chk_user_session_time_range} & 会话过期时间必须晚于创建时间，撤销时间不能早于创建时间。 \\
\code{course_schedule} & \code{chk_schedule_section_range} & 节次范围限定为 1 到 12，并要求结束节次不早于开始节次。 \\
\code{enrollment} & \code{UNIQUE(student_id, offering_id)} & 防止同一学生重复选择同一教学班。 \\
\end{longtable}

教学班容量与教室容量的关系属于跨表约束，普通 CHECK 无法查询另一张表，因此本系统使用触发器实现。迁移脚本执行前先审计并修正了已有样例数据中 \code{MA101} 班次容量大于 \code{M201} 教室容量的问题。

\section{跨表业务完整性与触发器设计}

本轮已经实现数据库触发器，触发器定义集中在 \code{opengauss_setup/sql/migrate_triggers_constraints_20260608.sql}。这些触发器不是用于维护 \code{selected_count}，而是用于角色档案一致性、先修关系合法性、选课完整性、教室容量一致性和成绩审计。表 \ref{tab:trigger-list} 给出触发器清单。

{\scriptsize
\setlength{\tabcolsep}{2pt}
\renewcommand{\arraystretch}{1.16}
\begin{longtable}{p{0.17\textwidth}p{0.12\textwidth}p{0.16\textwidth}p{0.25\textwidth}p{0.16\textwidth}p{0.08\textwidth}}
\caption{数据库触发器实现清单}\label{tab:trigger-list}\\
\toprule
触发器名称 & 作用表 & 触发时机 & 检查或维护内容 & 对应课程知识点 & 实现状态 \\
\midrule
\endfirsthead
\toprule
触发器名称 & 作用表 & 触发时机 & 检查或维护内容 & 对应课程知识点 & 实现状态 \\
\midrule
\endhead
\bottomrule
\endfoot
\code{trg_student_role_check} & \code{student} & 插入或修改 \code{user_id} 前 & 学生档案必须引用 \code{role='student'} 的账户 & 跨表用户自定义完整性 & 已实现 \\
\code{trg_teacher_role_check} & \code{teacher} & 插入或修改 \code{user_id} 前 & 教师档案必须引用 \code{role='teacher'} 的账户 & 子类型一致性 & 已实现 \\
\code{trg_admin_role_check} & \code{admin_profile} & 插入或修改 \code{user_id} 前 & 管理员档案必须引用 \code{role='admin'} 的账户 & 角色完整性 & 已实现 \\
\code{trg_prerequisite_no_cycle} & \code{course_prerequisite} & 插入或更新前 & 禁止自指先修课，并用递归查询检测环路 & 递归查询与复杂完整性 & 已实现 \\
\code{trg_offering_capacity_check} & \code{course_offering} & 插入或改容量、教室前 & 教学班容量不能超过教室容量，且不能低于当前已选人数 & 跨表约束 & 已实现 \\
\code{trg_classroom_capacity_update_check} & \code{classroom} & 更新 \code{capacity} 前 & 降低教室容量时不能小于已有教学班容量 & 更新约束 & 已实现 \\
\code{trg_enrollment_guard} & \code{enrollment} & 插入或改学生、班次、状态前 & 检查学生状态、班次状态、选课窗口、容量、时间冲突、先修课和退课限制 & 主动规则、事务与锁 & 已实现 \\
\code{trg_score_change_log} & \code{enrollment} & 更新 \code{final_score} 后 & 自动写入 \code{score_change_log}，记录旧成绩、新成绩、操作者和原因 & 审计触发器 & 已实现 \\
\end{longtable}
}

\section{选课完整性触发器}

\code{trg_enrollment_guard} 是选课系统最核心的数据库层完整性保护。对 \code{status='selected'} 的插入或恢复选课，它会锁定目标 \code{course_offering} 行，然后检查当前已选人数是否小于 \code{max_capacity}。已选人数通过 \code{enrollment} 聚合统计得到，系统没有恢复 \code{course_offering.selected_count}。

该触发器还基于 \code{course_schedule} 检查同一学期内的时间冲突，区间重叠条件为：

\begin{lstlisting}[style=sqlstyle]
target.weekday = existing.weekday
AND existing.start_section <= target.end_section
AND existing.end_section >= target.start_section
\end{lstlisting}

先修课程检查基于 \code{course_prerequisite} 和历史 \code{enrollment} 记录。系统采用 60 分作为通过先修课程的业务规则，即学生必须存在 \code{status='completed'} 且 \code{final_score >= 60} 的先修课程记录。退课时，触发器禁止已完成课程或已有成绩记录改为 \code{dropped}。

\section{成绩变更审计触发器}

\code{trg_score_change_log} 在 \code{enrollment.final_score} 更新后触发。当旧成绩和新成绩不同，触发器自动向 \code{score_change_log} 插入日志。由于日志表要求 \code{changed_by_user_id} 非空，业务层必须在同一事务内设置当前操作者：

\begin{lstlisting}[style=sqlstyle]
SELECT set_config('app.current_user_id', '<user_id>', true);
SELECT set_config('app.score_change_reason', '<reason>', true);
UPDATE enrollment
SET final_score = 88, status = 'completed'
WHERE enrollment_id = 1;
\end{lstlisting}

如果业务层没有设置 \code{app.current_user_id}，触发器会拒绝成绩更新。这比单纯在应用层手工插入日志更可靠，因为任何直接更新 \code{final_score} 的数据库操作都会被同一个审计规则约束。

\section{角色档案一致性触发器}

\code{student}、\code{teacher} 和 \code{admin_profile} 虽然都通过 \code{user_id} 外键引用 \code{user_account}，但普通外键只能保证账户存在，不能保证角色匹配。因此本系统实现了通用函数 \code{check_user_role_fn(p_user_id, p_expected_role)}，并由三个档案触发器复用。测试脚本已经验证，将教师账户插入学生档案、将学生账户插入教师档案都会被拒绝。

\section{先修课程合法性触发器}

\code{chk_prereq_not_self} 可以阻止 \code{course_id = prereq_course_id} 的直接自指，但不能发现更复杂的 A 依赖 B、B 依赖 C、C 又依赖 A 的环路。本系统通过 \code{trg_prerequisite_no_cycle} 在插入或更新先修关系前执行递归查询，沿着先修链查找是否会回到当前课程。如果形成环路，触发器抛出异常并阻止写入。

\section{应用层校验与数据库触发器的关系}

当前应用层仍保留友好的错误提示和页面交互检查，例如 \code{services/selection_service.py} 会调用数据库函数 \code{select_course_tx()}，\code{services/score_service.py} 会在事务内设置审计上下文。数据库触发器承担兜底一致性责任：即使绕过前端或服务层，只要直接写数据库，也必须通过容量、时间冲突、先修课程、角色档案和成绩审计规则。

这种分层体现了数据库完整性设计思想：局部、稳定的规则用 CHECK、UNIQUE 和外键表达；跨表和上下文相关规则用触发器表达；用户体验相关的提示和权限入口仍由应用层负责。

\section{本章小结}

本章说明了当前系统已经从“约束设计方案”推进到“数据库层实现”。主键、外键、唯一约束、CHECK、DEFAULT 保证基础完整性；角色、容量、先修、选课和成绩日志等跨表规则由触发器维护；应用层负责权限入口、事务调用和友好错误提示。尤其需要强调的是，触发器没有恢复 \code{selected_count}，而是在保留 3NF 结构的前提下承担复杂完整性检查和审计职责。
